\section{Related Work}

\subsection{Outcome-Based Reinforcement Learning for Reasoning}

Reinforcement learning from outcome feedback is widely used to align language models and improve their reasoning behavior. PPO-style RLHF optimizes a policy against learned or programmatic rewards while controlling policy drift \cite{christiano2017deep,schulman2017proximal,ouyang2022training}. Direct preference optimization and related preference-learning methods simplify parts of this pipeline when pairwise preference data are available \cite{rafailov2023direct}. For mathematical reasoning, recent systems often use verifiable rewards, where a final answer can be checked automatically \cite{shao2024deepseekmath,guo2025deepseekr1}. GRPO and related group-relative methods reduce the need for a learned value function by comparing multiple responses sampled for the same prompt and broadcasting a scalar relative advantage over response tokens \cite{shao2024deepseekmath}. These methods use reliable final feedback, but they do not provide fine-grained token-level guidance. POV-OPD addresses this limitation by retaining dense teacher supervision while using final correctness as a validation signal.

\subsection{Knowledge Distillation and On-Policy Distillation}

Knowledge distillation transfers behavior from a stronger teacher model to a smaller or weaker student \cite{hinton2015distilling}. For sequence generation, distillation can imitate teacher outputs or teacher distributions at the token or sequence level. Offline distillation is simple, but it trains the student on states induced by the teacher or by a fixed dataset, which can differ from the states visited by the student during reasoning. On-policy distillation reduces this mismatch by evaluating teacher preferences on student-generated trajectories. Recent work on OPD shows that dense token-level teacher rewards can be useful but sensitive to teacher-student compatibility and the states visited by the student \cite{li2026rethinkingopd}. POV-OPD follows the same on-policy view: the student samples the response, the teacher scores the visited prefixes, and the student updates from dense teacher-student log-probability gaps. The key difference is that POV-OPD does not assume every dense teacher preference is equally reliable.

\subsection{Credit Assignment and Dense Supervision}

Long-horizon reasoning requires assigning credit to intermediate decisions from sparse final outcomes. Chain-of-thought prompting, self-consistency, and self-training methods highlight the importance of intermediate reasoning paths for final-answer accuracy \cite{wei2022chain,wang2022selfconsistency,zelikman2022star}. Process supervision provides dense feedback for intermediate steps, but such supervision can be costly or noisy \cite{lightman2023lets}. Dense teacher rewards offer an automatic alternative, yet they can conflict with the final correctness of the trajectory. POV-OPD treats this conflict as a validation problem: final correctness determines whether the local teacher-student preference points in an outcome-consistent direction. Our prefix validation is also related to classical change detection, where CUSUM statistics are used to detect sustained shifts in a monitored signal \cite{page1954continuous,basseville1993detection}. Here the monitored signal is teacher excess surprisal along a student-generated prefix, and the response is a monotone decay of later OPD weights rather than a hard rejection of the trajectory.
